\section{Requisiti funzionali e non funzionali}

\subsection{Requisiti funzionali}
\begin{longtable}{@{}p{1.2cm}p{4.5cm}X@{}}
    \toprule
    \textbf{ID} & \textbf{Nome} & \textbf{Descrizione} \\
    \midrule
    \endfirsthead
    \toprule
    \textbf{ID} & \textbf{Nome} & \textbf{Descrizione} \\
    \midrule
    \endhead

    \FRid{1} & Ricerca per signature &
    Il sistema deve consentire di cercare un testo per
    \code{track\_name}~+~\code{artist\_name}, con eventuale filtro su
    \code{album\_name} e \code{duration}. Realizzato da
    \endpoint{GET}{/api/get} e \endpoint{GET}{/api/get-cached}. \\

    \FRid{2} & Recupero per ID &
    Deve essere possibile recuperare un testo conoscendo il suo identificativo
    univoco. Realizzato da \endpoint{GET}{/api/get/<id>}. \\

    \FRid{3} & Ricerca generica &
    L'utente deve poter cercare testi per parola chiave (\code{q}) o per
    combinazione di titolo, artista, album. Realizzato da
    \endpoint{GET}{/api/search}. \\

    \FRid{4} & Pubblicazione anonima &
    Un utente non registrato deve poter pubblicare un nuovo testo dopo aver
    superato una sfida \emph{Proof of Work}. Realizzato dalla coppia
    \endpoint{POST}{/api/request-challenge} +
    \endpoint{POST}{/api/publish}~con header \code{X-Publish-Token}. \\

    \FRid{5} & Pubblicazione autenticata &
    Un utente autenticato (Bearer token) deve poter pubblicare un nuovo testo
    senza risolvere la PoW. Il brano viene attribuito al suo username nel
    campo \code{submittedBy}. \\

    \FRid{6} & Registrazione utente &
    Deve essere possibile creare un account fornendo username, email e
    password. Realizzato da \endpoint{POST}{/auth/register}. \\

    \FRid{7} & Login utente &
    L'utente deve potersi autenticare con username o email e password,
    ricevendo \emph{access} e \emph{refresh token}. Realizzato da
    \endpoint{POST}{/auth/login}. \\

    \FRid{8} & Refresh del token &
    L'\emph{access token} scaduto deve poter essere rinnovato senza re-login,
    presentando il \emph{refresh token}. Realizzato da
    \endpoint{POST}{/auth/refresh}. \\

    \FRid{9} & Dati dell'utente corrente &
    L'utente autenticato deve poter recuperare i propri dati di profilo.
    Realizzato da \endpoint{GET}{/auth/me}. \\

    \FRid{10} & Elenco dei propri brani &
    L'utente autenticato deve poter ottenere la lista dei testi da lui
    pubblicati. Realizzato da \endpoint{GET}{/api/me/songs}. \\

    \FRid{11} & Modifica di un proprio brano &
    L'utente autenticato deve poter modificare un brano di cui è autore.
    Realizzato da \endpoint{PUT}{/api/songs/<id>}. \\

    \FRid{12} & Cancellazione di un proprio brano &
    L'utente autenticato deve poter cancellare un brano di cui è autore.
    Realizzato da \endpoint{DELETE}{/api/songs/<id>}. \\

    \FRid{13} & Esplorazione del catalogo &
    Il sistema deve fornire un elenco paginato dei brani, ordinabile per
    data, titolo o artista. Realizzato da \endpoint{GET}{/explore}. \\

    \FRid{14} & Lettura del testo sincronizzato &
    L'applicazione mobile deve poter mostrare il testo sincronizzato a tempo
    con la riproduzione, evidenziando la riga corrente. \\

    \FRid{15} & Health check &
    Deve essere disponibile un endpoint per verificare lo stato del backend.
    Realizzato da \endpoint{GET}{/health}. \\

    \bottomrule
    \caption{Requisiti funzionali.}
\end{longtable}

\subsection{Requisiti non funzionali}
\begin{longtable}{@{}p{1.5cm}p{2.6cm}p{6cm}p{3cm}@{}}
    \toprule
    \textbf{ID} & \textbf{Tipo} & \textbf{Descrizione} & \textbf{Riferito a} \\
    \midrule
    \endfirsthead
    \toprule
    \textbf{ID} & \textbf{Tipo} & \textbf{Descrizione} & \textbf{Riferito a} \\
    \midrule
    \endhead

    \NFRid{1} & Usabilità &
    L'app mobile deve essere comprensibile da un utente non tecnico in meno
    di un minuto di navigazione.
    & Mobile app \\

    \NFRid{2} & Prestazioni &
    Le risposte dell'API devono avere latenza p95 $<$ 200~ms su catalogo
    fino a 10000 brani in ambiente di sviluppo.
    & Backend \\

    \NFRid{3} & Portabilità &
    Il codice del backend deve girare senza modifiche su Linux e macOS;
    l'app mobile su iOS, Android e Web.
    & Tutto il sistema \\

    \NFRid{4} & Manutenibilità &
    Architettura a strati \emph{routes/services/models/utils} con
    separazione netta delle responsabilità.
    & Backend \\

    \NFRid{5} & Sicurezza --- riservatezza &
    Le password degli utenti non devono mai essere memorizzate in chiaro
    né essere visibili nei log.
    & DB, Auth \\

    \NFRid{6} & Sicurezza --- integrità &
    Una pubblicazione non può essere attribuita all'utente sbagliato; un
    utente non può modificare o cancellare brani di cui non è autore.
    & API \\

    \NFRid{7} & Sicurezza --- disponibilità &
    Il sistema deve rifiutare richieste con body oltre la soglia
    (\code{MAX\_CONTENT\_LENGTH = 256 KB}) per limitare attacchi DoS.
    & Backend \\

    \NFRid{8} & Scalabilità &
    Il backend deve poter girare dietro a un WSGI server multi-worker
    (es. \emph{gunicorn}) senza necessità di modifiche significative.
    & Backend \\

    \NFRid{9} & Conformità --- GDPR &
    L'utente deve poter accedere ai propri dati personali
    (\code{/auth/me}). La cancellazione dell'account è prevista come
    estensione futura.
    & Auth \\

    \NFRid{10} & Documentabilità &
    Il sistema deve essere accompagnato da un \code{README} e da una
    documentazione delle API consultabile via esempi \code{curl}.
    & Repository \\

    \bottomrule
    \caption{Requisiti non funzionali.}
\end{longtable}

\subsection{Requisiti di sicurezza}
\begin{longtable}{@{}p{1.5cm}p{4cm}X@{}}
    \toprule
    \textbf{ID} & \textbf{Categoria} & \textbf{Descrizione} \\
    \midrule
    \endfirsthead
    \toprule
    \textbf{ID} & \textbf{Categoria} & \textbf{Descrizione} \\
    \midrule
    \endhead

    SEC1 & Autenticazione &
    Username/email + password con hashing PBKDF2-SHA256 (\emph{salt}
    randomico per record). Le risposte di login non distinguono fra
    utente inesistente e password errata. \\

    SEC2 & Autorizzazione &
    Le route protette verificano il JWT tramite \code{@jwt\_required}.
    Le mutazioni su brani richiedono che
    \code{song.user\_id == current\_user\_id}. \\

    SEC3 & Anti-spam &
    La pubblicazione anonima richiede il superamento di una sfida PoW
    one-shot, marcata atomicamente come consumata via UPDATE WHERE. \\

    SEC4 & Validazione input &
    Validazione lato server con espressioni regolari su username
    (\verb|^[A-Za-z0-9_]{3,30}$|), email (regex permissiva), password
    ($\geq 8$ caratteri, almeno una cifra) e limiti di lunghezza per i
    campi testuali. \\

    SEC5 & Protezione DoS &
    Limite del body della richiesta a 256~KB, limiti per campo
    (50~KB per \code{plainLyrics}, 100~KB per \code{syncedLyrics},
    300~caratteri per i campi anagrafici). \\

    SEC6 & Cifratura in transito (deferred) &
    In produzione: terminazione TLS a livello di reverse proxy con header
    HSTS. \\

    SEC7 & Cifratura at-rest (deferred) &
    In produzione: cifratura del filesystem dove risiede \code{lyrics.db}. \\

    SEC8 & Logging (deferred) &
    Log strutturato dei login falliti e delle azioni di pubblicazione,
    con rotazione su file. \\

    SEC9 & Backup (deferred) &
    Snapshot periodico del file SQLite con retention $\geq 7$ giorni. \\

    SEC10 & Conformità --- GDPR &
    Trattamento dei dati personali documentato; endpoint
    \code{DELETE /auth/me} previsto come estensione futura per il
    diritto all'oblio. \\

    \bottomrule
    \caption{Requisiti di sicurezza.}
\end{longtable}
